\documentclass[../DoAn.tex]{subfiles}
\begin{document}
\section{Đặt vấn đề}
\label{section:1.1}

Hoạt động tư vấn pháp luật đóng vai trò cốt lõi trong việc cung cấp hỗ trợ pháp lý, giải đáp các vướng mắc và hướng dẫn người dân tuân thủ đúng quy định, nhằm bảo vệ tối đa quyền và lợi ích hợp pháp của các cá nhân cũng như tổ chức.

Trong hệ thống pháp lý, Luật Hôn nhân và Gia đình là một lĩnh vực chuyên ngành đặc thù, điều chỉnh trực tiếp các mối quan hệ phát sinh từ hôn nhân và gia đình. Phân hệ pháp luật này bao gồm toàn bộ các quy định chặt chẽ về việc kết hôn, ly hôn, quyền và nghĩa vụ của vợ chồng, quyền nuôi con, phân chia khối tài sản chung và các vấn đề khác liên quan đến gia đình. Mục tiêu tối thượng của Luật Hôn nhân và Gia đình là bảo vệ quyền lợi hợp pháp của các bên trong mối quan hệ hôn nhân và gia đình, đồng thời đảm bảo sự công bằng và ổn định trong xã hội.

Hiện nay, nhu cầu nhận tư vấn chuyên sâu về lĩnh vực hôn nhân và gia đình đang tăng trưởng mạnh mẽ do các tác động xã hội phức tạp, cụ thể bao gồm các nguyên nhân như tỷ lệ ly hôn tăng và trẻ hóa, tranh chấp tài sản và quyền nuôi con, xu hướng lập thỏa thuận tài sản và các vấn đề hôn nhân khác như kết hôn có yếu tố nước ngoài, mang thai hộ... 

Trước nhu cầu cấp thiết đó, việc nghiên cứu và xây dựng một hệ thống chatbot có khả năng tư vấn tự động về Luật Hôn nhân và Gia đình mang lại những ý nghĩa thực tiễn vô cùng to lớn. Hệ thống này đóng vai trò như một công cụ bổ trợ đắc lực cho các cơ quan tư vấn công lập hiện tại, giúp giảm tải hiệu quả áp lực công việc cho đội ngũ cán bộ ngành luật. Đồng thời, giải pháp này giúp người dân đại chúng nhanh chóng tiếp cận được những tư vấn thiết thực, hợp pháp và hoàn toàn sát hợp với bối cảnh, hoàn cảnh cá nhân của riêng mình.

Mặc dù mang ý nghĩa quan trọng, việc hiện thực hóa một chatbot tư vấn pháp lý tự động vẫn phải đối mặt với những rào cản mang tính đặc thù. Hệ thống pháp luật về hôn nhân và gia đình tại Việt Nam mang tính chất phức tạp, gồm nhiều lớp văn bản (Luật, Nghị định, Thông tư, Nghị quyết…) và thường xuyên được cập nhật, gây khó khăn trong tiếp cận, tra cứu và áp dụng. Một khó khăn quan trọng khác là các văn bản pháp luật có thể tồn tại quan hệ chồng chéo, sửa đổi, bổ sung, thay thế, hướng dẫn hoặc thậm chí có khả năng mâu thuẫn trong một số tình huống áp dụng cụ thể. Người dùng không chuyên thường khó xác định văn bản nào còn hiệu lực, văn bản nào là căn cứ chính, văn bản nào chỉ mang tính hướng dẫn hoặc bổ trợ, và văn bản nào không nên áp dụng do đã bị thay thế hoặc không phù hợp với phạm vi tình huống. Vì vậy, một chatbot tư vấn pháp luật không chỉ cần truy xuất đúng điều luật, mà còn cần kiểm tra quan hệ giữa các văn bản viện dẫn, xử lý khả năng xung đột căn cứ pháp lý và giải thích rõ vì sao các căn cứ đó được sử dụng trong câu trả lời.

Hiện nay, các công cụ hỗ trợ tra cứu và tư vấn pháp luật tại Việt Nam đã có những bước phát triển nhất định với sự xuất hiện của các sản phẩm như AITracuuluat, AI luật và Chatbot của Bộ Tư pháp. Bên cạnh đó, các mô hình ngôn ngữ lớn phổ biến như Gemini cũng thường xuyên được người dùng sử dụng để giải đáp các vướng mắc pháp lý. Trong lĩnh vực chuyên biệt, các doanh nghiệp viễn thông lớn đã triển khai các hệ thống như Trợ lý ảo Viettel hay VNPT SmartBot phục vụ cho các tổ chức và cá nhân có thẩm quyền.

Qua quá trình khảo sát và đánh giá, các sản phẩm hiện có bộc lộ những ưu điểm và hạn chế khác nhau trong việc giải quyết các bài toán pháp lý. Công cụ AITracuuluat đôi khi nêu ra các căn cứ pháp lý không chính xác, không liên quan đến tình huống trong câu hỏi của người dùng hoặc gặp tình trạng dữ liệu điều luật chưa được cập nhật mới nhất. Trong khi đó, chatbot của Bộ Tư Pháp mặc dù đảm bảo tính cập nhật và độ chính xác cao , hệ thống này lại chưa có khả năng cảnh báo các văn bản đã hoặc sắp hết hiệu lực, chưa giải thích lí do vì sao áp dụng các căn cứ pháp lý được nêu ra trong câu trả lời và bị giới hạn về số lượng lượt sử dụng miễn phí Đối với các mô hình ngôn ngữ lớn như Gemini: Câu trả lời thường xuyên thiếu sự viện dẫn đầy đủ các căn cứ pháp lý cần thiết và gặp khó khăn trong việc cập nhật thời hạn hiệu lực của văn bản quy phạm pháp luật. Các hệ thống chuyên dụng của Viettel hay VNPT thì khả năng tiếp cận của người dân đại chúng còn hạn chế do quy định về quyền truy cập.

Tựu trung lại, các chatbot tư vấn pháp luật Việt Nam hiện hành còn gặp phải 3 tồn tại tiêu biểu: (i) Dữ liệu điều luật căn cứ chưa được cập nhật kịp thời, đầy đủ hoặc đôi khi xảy ra hiện tượng trích dẫn thiếu chính xác, không sát thực tế tình huống được hỏi. (ii) Chưa có cơ chế cảnh báo tình trạng hiệu lực văn bản (chưa cảnh báo văn bản đã/sắp hết hiệu lực) và giải thích lý do vì sao áp dụng các căn cứ đó. (iii) Khả năng tiếp cận của người dân bị thu hẹp đáng kể do các rào cản về quyền truy cập nội bộ hoặc bị giới hạn về số lượng câu hỏi sử dụng miễn phí.

Việc giải quyết bài toán tra cứu và tư vấn pháp luật tự động mang lại những giá trị thực tiễn thiết thực cho cộng đồng. Xử lý những rào cản thông tin nêu trên giúp người dân nhanh chóng tiếp cận nguồn dữ liệu pháp lý minh bạch, chính thống và sát với hoàn cảnh cá nhân. Một cơ chế tư vấn luật pháp hiệu quả sẽ đóng vai trò bổ trợ đắc lực cho các hệ thống tư vấn công lập, góp phần giảm thiểu áp lực công việc cho đội ngũ cán bộ ngành luật. Khuôn mẫu giải quyết bài toán này không chỉ giới hạn trong phạm vi Luật Hôn nhân và gia đình mà còn sở hữu tiềm năng mở rộng sang nhiều phân hệ pháp luật chuyên ngành khác.
\section{Mục tiêu và phạm vi đề tài}
\label{section:1.2}
Từ những đánh giá, phân tích về các vấn đề, tồn tại trên, ĐATN hướng tới mục tiêu xây dựng một hệ thống chatbot tư vấn Luật Hôn nhân và gia đình có khả năng đưa ra những câu trả lời tham chiếu đúng các điều luật liên quan, cảnh báo người dùng về hiệu lực của các căn cứ pháp lý, giải thích ngắn gọn quan hệ giữa các văn bản viện dẫn và giải thích vì sao áp dụng các căn cứ pháp lý đó. Trọng tâm của đồ án là thiết lập một cơ chế truy xuất thông tin đảm bảo mọi phản hổi đều dựa trên các các căn cứ pháp lý còn hiệu lực pháp luật tại thời điểm truy vấn, các căn cứ pháp lý được áp dụng theo đúng nguyên tắc hiệu lực các văn bản quy phạm pháp luật của hệ thống luật pháp Việt Nam và có thể xử lý các vấn đề chồng chéo, mâu thuẫn giữa các điều luật (nếu có).

Phạm vi của đề tài là phần mềm được phát triển có khả năng trả lời, giải quyết các câu hỏi thuộc các chủ đề liên quan đến luật hôn nhân gia đình, bao gồm: Kết hôn, Quan hệ giữa vợ và chồng trong hôn nhân, Chấm dứt hôn nhân, Chế độ tài sản của vợ chồng, Quan hệ giữa cha mẹ và con, Quan hệ hôn nhân và gia đình có yếu tố nước ngoài. Phần mềm tập trung vào các chức năng chính bao gồm: hỏi đáp về nội dung Luật Hôn nhân và gia đình Việt Nam; phân loại dạng câu hỏi theo các chủ đề pháp lý và truy xuất ngữ cảnh pháp lý từ đồ thị tri thức (Knowledge Graph) nhờ các agent; và sinh câu trả lời có trích dẫn điều khoản cụ thể. Hệ thống cung cấp khả năng hiển thị minh bạch luồng xử lý của các tác nhân (agent) và các công cụ được kích hoạt trong quá trình tìm kiếm câu trả lời. Công cụ cũng tích hợp cơ chế kiểm tra và cảnh báo về tình trạng hiệu lực, thứ bậc văn bản và các mối quan hệ sửa đổi, bổ sung của các căn cứ pháp lý được sử dụng. Việc thực hiện đề tài kỳ vọng đạt được độ chính xác trong viện dẫn pháp lý và khả năng xử lý xung đột văn bản ở mức trên 85\%, góp phần nâng cao chất lượng dịch vụ tư vấn pháp luật tự động cho cộng đồng.

\section{Định hướng giải pháp}
\label{section:1.3}

ĐATN đã khảo sát, phân tích, nêu ra vấn đề của các ứng dụng chatbot tư vấn pháp luật hiện hành tại Việt Nam cùng những hạn chế liên quan đến việc trích dẫn các căn cứ pháp lý chính xác là cơ sở để trả lời cho câu hỏi của người dùng, thiếu cơ chế cảnh báo hiệu lực văn bản và xử lý các tình huống mâu thuẫn pháp lý thực tế. Để đáp ứng, hoàn thiện, khắc phục những hạn chế, tồn tại đã chỉ ra ở trên, ĐATN đưa ra phương án giải quyết cho từng vấn đề: (i) mô hình hóa các thực thể pháp lý và mối liên hệ giữa chúng (ví dụ: giữa “quyền nuôi con” và “độ tuổi trẻ em” hay “thỏa thuận của cha mẹ”) và các mối quan hệ hiệu lực, sửa đổi, hướng dẫn bổ sung, thay thế, bác bỏ giữa các văn bản quy phạm pháp luật của Luật Hôn nhân và Gia đình Việt Nam dưới dạng đồ thị tri thức (Knowledge Graph); (ii) Agentic RAG để tối ưu việc truy vấn trên đồ thị tri thức bằng cách phân tách các retriever tương ứng với các dạng câu hỏi và sử dụng agent để quyết định sẽ sử dụng (những) retriever nào để lấy về dữ liệu phục vụ cho việc sinh câu trả lời. và (iii) thiết kế cơ chế tự động kiểm tra và cảnh báo tình trạng hiệu lực, đồng thời phát hiện và xử lý mâu thuẫn trong các căn cứ pháp lý (nếu có) để sinh câu trả lời chính xác, minh bạch.

Từ đầu vào là một câu hỏi của người dùng bằng tiếng Việt tự nhiên thuộc các chủ đề của Luật Hôn nhân và Gia đình , hệ thống đề xuất cần đưa ra đầu ra bao gồm một câu trả lời hoàn chỉnh đính kèm trích dẫn điều khoản chính xác, hiển thị tường minh luồng xử lý của các agent (agent trace) và ngữ cảnh pháp lý được truy xuất từ đồ thị tri thức. Trước hết, kết quả phản hồi này cần thỏa mãn ràng buộc về tính chính xác pháp lý tuyệt đối, đảm bảo các điều luật được viện dẫn đều là văn bản đang còn hiệu lực hoặc hoàn toàn phù hợp với mốc thời gian và bối cảnh của tình huống được hỏi. Ngoài ra, trong câu trả lời cần có những cảnh báo pháp lý khi áp dụng các căn cứ pháp lý đã hết hiệu lực hoặc sắp hết hiệu lực, giải thích tại sao lại áp dụng các căn cứ pháp lý đó trong trường hợp có mâu thuẫn, chồng chéo trong những căn cứ pháp lý.

Quy trình xử lý của kiến trúc đề xuất bao gồm 4 bước chính, cụ thể thực hiện theo thứ tự như sau. Ban đầu, sử dụng mô hình ngôn ngữ lớn kết hợp Prompt Engineering để phân loại dạng câu hỏi và xác định chủ đề pháp lý liên quan nhằm thu hẹp phạm vi tìm kiếm. Tiếp theo, điều phối các Agent Retriever chuyên để xử lý cho từng dạng câu hỏi, các agent có nhiệm vụ phải trích xuất ra các ý định (intent) và thực thể (entity) trong câu hỏi của người dùng, sau đó điền các tham số trích xuất được vào mẫu câu lệnh cypher (cypher template) để sinh ra câu lệnh Cypher để truy xuất các căn cứ pháp lý còn hiệu lực (hoặc phù hợp với thời điểm xảy ra sự kiện pháp lý trong câu hỏi) từ cơ sở dữ liệu đồ thị Neo4j. Tiếp theo, áp dụng cơ chế lọc và sắp xếp căn cứ pháp lý theo thứ bậc hiệu lực và thời điểm ban hành để phát hiện cũng như xử lý các trường hợp chồng chéo hoặc mâu thuẫn (nếu có). Cuối cùng, gọi API mô hình ngôn ngữ lớn để sinh câu trả lời cuối cùng đảm bảo có trích dẫn căn cứ pháp lý chi tiết, cảnh báo hiệu lực văn bản và giải thích ngắn gọn mối quan hệ giữa các căn cứ pháp lý, lí do chọn các căn cứ pháp lý đó. 

Công nghệ được sử dụng cho ĐATN bao gồm nhiều thành phần khác nhau. Thư viện Chainlit hỗ trợ việc lập trình giao diện Web Chatbot UI một cách nhanh chóng, cung cấp sẵn các cơ chế quản lý session hỏi đáp và hiển thị trực quan luồng agent trace cho người dùng. Ngôn ngữ lập trình Python được chọn làm ngôn ngữ chủ đạo nhờ sở hữu hệ sinh thái thư viện phong phú, tối ưu cho các bài toán xử lý ngôn ngữ tự nhiên và khoa học dữ liệu. Cơ sở dữ liệu đồ thị Neo4j cho phép mô hình hóa, lưu trữ và truy vấn nhanh chóng các mối quan hệ thực thể pháp lý, kết hợp với PostgreSQL nhằm quản lý thông tin người dùng và lưu trữ lịch sử hội thoại. Để thực hiện các tác vụ suy luận và trích xuất dữ liệu có cấu trúc, hệ thống tích hợp API của các mô hình ngôn ngữ lớn tiên tiến bao gồm dòng GPT-4.1, gpt 3o và gpt-4o, Gemini 3.1  nhằm đảm bảo chất lượng sinh ngôn ngữ. Cuối cùng, thư viện RAGAS được ứng dụng làm nền tảng đánh giá định lượng, giúp đo lường chính xác hiệu năng của toàn bộ luồng truy xuất.

Các đóng góp chính của đồ án bao gồm: (i) đề xuất và xây dựng thành công kiến trúc Agentic GraphRAG truy xuất thông tin trên đồ thị tri thức phục vụ tư vấn pháp luật;  (ii) thiết kế cấu trúc Schema Đồ thị tri thức toàn diện biểu diễn rõ ràng các quan hệ hiệu lực, hướng dẫn, sửa đổi và xung đột giữa các văn bản pháp luật thuộc Luật Hôn nhân và Gia đình; (iii) xây dựng cơ chế tự động phát hiện, cảnh báo và ưu tiên áp dụng căn cứ pháp lý khi xuất hiện tình trạng chồng chéo hoặc mâu thuẫn văn bản (iv) thiết kế và xây dựng bộ Benchmark dataset tiêu chuẩn gồm tối thiểu 600 câu hỏi - đáp pháp lý có gán nhãn cấu trúc chi tiết phục vụ kiểm thử; (v) phát triển hoàn chỉnh ứng dụng Web Chatbot trên nền tảng Chainlit tích hợp đầy đủ tính năng hiển thị luồng agent trace và tài liệu trích dẫn.

\section{Bố cục đồ án}
\label{section:1.4}

Phần còn lại của báo cáo đồ án tốt nghiệp này được tổ chức như sau.

Chương 2 trình bày tổng quan kết quả khảo sát các hệ thống chatbot tư vấn pháp luật hiện hành và thực trạng áp dụng trợ lý ảo để tư vấn luật pháp nói chung và luật hôn nhân gia đình nói riêng tại Việt Nam. Căn cứ vào các ưu nhược điểm của những công cụ đang có trên thị trường, các giải pháp hiện hành và lấy đó làm cơ sở để xác định rõ bài toán mà đồ án cần tập trung giải quyết.

Trong Chương 3, đồ án giới thiệu các cơ sở lý thuyết và nền tảng công nghệ lõi được ứng dụng để phát triển hệ thống phần mềm. Nội dung tập trung làm rõ các khái niệm về Đồ thị tri thức (Knowledge Graph), kỹ thuật Retrieval-Augmented Generation (RAG) và kiến trúc Agentic AI. Bên cạnh đó, chương này cũng đề cập đến các công cụ lưu trữ dữ liệu đồ thị như Neo4j và nền tảng giao diện Chainlit.

Chương 4 cung cấp nội dung chi tiết về quá trình thiết kế hệ thống và quá trình thử nghiệm hệ thống và phân tích các kết quả đánh giá hiệu năng. Dựa trên bộ dữ liệu kiểm chuẩn chuẩn hóa (benchmark dataset), chương này trình bày các số liệu thực nghiệm liên quan đến độ chính xác trong việc lựa chọn công cụ của tác nhân AI, khả năng truy xuất ngữ cảnh và chất lượng câu trả lời cuối cùng. Các kết quả này được đối chiếu trực tiếp với hệ thống chỉ tiêu đề ra thông qua bộ chỉ số RAGAS và các tiêu chí xử lý văn bản luật khác.

Chương 5 trình bày đóng góp cốt lõi của đồ án thông qua việc mô tả chi tiết giải pháp kỹ thuật và quá trình xây dựng hệ thống Agentic GraphRAG. Chương này mô tả việc thiết kế đồ thị knowledge graph để mô hình hóa các thực thể pháp lý cùng mối quan hệ giữa Luật hôn nhân và gia đình năm 2014 và các văn bản dưới luật liên qua và cơ chế tự động kiểm tra tình trạng hiệu lực văn bản, áp dụng các văn bản theo đúng cấp bậc các văn bản pháp luật và giải quyết mâu thuẫn giữa các điều luật (nếu có). Cuối cùng, trình bày cách xây dựng tập benchmark dataset gồm 600 câu hỏi, bao phủ các dạng câu hỏi chính liên quan đến Luật hôn nhân và gia đình.

Cuối cùng, Chương 6 tổng kết lại các kết quả đã đạt được của quá trình nghiên cứu và phát triển phần mềm. Bên cạnh việc khẳng định mức độ hoàn thành các mục tiêu ban đầu, chương này cũng đưa ra những nhận định khách quan về các mặt còn hạn chế của hệ thống hiện tại, từ đó đề xuất các hướng phát triển và mở rộng tiềm năng của công cụ tư vấn pháp lý trong tương lai.

\end{document}x